\documentclass[11pt]{article}

\usepackage[T1]{fontenc}
\usepackage[utf8]{inputenc}
\usepackage{lmodern}
\usepackage{microtype}
\usepackage[a4paper,margin=27mm,headheight=14pt]{geometry}
\usepackage{amsmath,amssymb,amsthm,amscd,mathtools,mathrsfs}
\usepackage{enumitem}
\usepackage{xcolor}
\usepackage{booktabs,array}
\usepackage{fancyhdr}
\usepackage[colorlinks=true,linkcolor=blue!45!black,
 citecolor=blue!45!black,urlcolor=blue!45!black]{hyperref}
\usepackage[nameinlink,noabbrev]{cleveref}

\newtheorem{theorem}{Theorem}[section]
\newtheorem{proposition}[theorem]{Proposition}
\newtheorem{lemma}[theorem]{Lemma}
\newtheorem{corollary}[theorem]{Corollary}
\theoremstyle{definition}
\newtheorem{definition}[theorem]{Definition}
\theoremstyle{remark}
\newtheorem{remark}[theorem]{Remark}

\numberwithin{equation}{section}
\setlist{itemsep=2pt,topsep=4pt}
\setlength{\parindent}{0pt}
\setlength{\parskip}{0.48em}
\setlength{\emergencystretch}{2em}
\pagestyle{fancy}
\fancyhf{}
\fancyhead[L]{The \(\mathrm{SU}(3)\) Gram and toroidal carriers}
\fancyhead[R]{Verification draft}
\fancyfoot[C]{\thepage}

\title{\bfseries The \(\mathrm{SU}(3)\) Gram Cone and the Rank-Six
Toroidal Lattice\\
\large Exact orbit spaces, Weyl-chamber incidence, and the obstruction to
a monodromy-compatible Gram map}
\author{Independent mathematical verification draft}
\date{4 September 2026}

\begin{document}
\maketitle

\begin{abstract}
Let \(\mathrm{SU}(3)\) act on
\(\operatorname{Mat}_{3\times2}(\mathbb C)\) by left multiplication.  The
first part proves that the Gram map induces a proper stratified
homeomorphism from the orbit space onto the positive-semidefinite
\(2\times2\) Hermitian cone, determines every stabilizer and orientation,
and identifies the compact quotient as \(B^4\), not \(S^4\).  The second
part treats the distinct rank-six toroidal lattice
\(\Lambda_3=H_1(T^6;\mathbb Z)=L_3\oplus\bar L_3\).  We compute its exact
diagonal \(S_3\)-action, integral symplectic monodromy, and the quotient
\(\Phi=P\oplus Q\).  Its kernel is a fixed rank-two lattice, while each
rank-two quotient is the \(A_2\) weight lattice rather than the root lattice;
the sequence splits as abelian groups but has no integral
\(S_3\)-equivariant splitting.  The six Kuhn chambers form the regular
\(S_3\)-set, not the two three-point orbits furnished by a basis of
\(\Lambda_3\).  An exact incidence complex and the spectral-floor evaluation
place the six chambers on six diagonal rank-two Gram-boundary points, but
the tempting extension from all six lattice directions has a non-invariant
kernel and therefore cannot intertwine monodromy.  Finally we prove that the
branch-exponent map restricts to the angular map \(P\); composing it with the
deck map \(Q\) is a domain error.  The separate free defect lattice
\(H^4(U;\mathbb Z)\cong(\mathbb Z^2)^3\) is not identified with
\(\Lambda_3\), despite the equality of ranks.
\end{abstract}

\medskip
\noindent\fcolorbox{black!55}{black!3}{%
\begin{minipage}{0.94\linewidth}
\textbf{Status and scope.}
This note separates two exact carriers.  The rank-two Gram quotient follows
from the displayed \(\mathrm{SU}(3)\) action.  The later toroidal calculation
uses the explicitly marked lattice \(\Lambda_3\), its Kuhn fan, and
\(\Phi=P\oplus Q\); it does not assert a holomorphic map from the toroidal
total space to a Gram or Jordan cone.  The representation sphere
\(S^{12}_{\mathrm{col}}\) is not a claimed complex manifold and is not
automatically a cusp-induced smash sphere.  The note constructs neither a
nonzero instanton sector on \(S^4\) nor a quantum Yang--Mills Hamiltonian.
The rank-three toroidal lattice and the regular-locus defect group are kept
as different typed objects throughout.
\end{minipage}}

\paragraph{Provenance.}
This standalone draft extracts the rigorously surviving Gram-quotient and
toroidal-lattice calculations from independent audits of proposed
continuations of the reconstructed \(S^6\) manuscript.  The orbit,
stabilizer, orientation, compactification, clutching, lattice, incidence, and
obstruction claims are proved below from their stated objects; no theorem
about the existence of the candidate complex structure on \(S^6\) is used.

\section{The higher-rung \texorpdfstring{\(\mathrm{SU}(3)\)}{SU(3)} Gram cone}
\label{sec:higher-rung-su3-gram-cone}

This section isolates the part of the proposed higher-rung construction that
is independent of the candidate complex threefold and proves it from the
underlying linear action.  The compact object used here is the representation
sphere
\[
 \operatorname{Mat}_{3\times2}(\mathbb C)^+,
\]
the one-point compactification of a twelve-real-dimensional vector space.
It must not be identified with a cusp-induced \(S^{12}\) construction unless
an additional based identification has first been chosen.  The proofs below
also distinguish three structures that otherwise look deceptively similar:
the Lorentzian determinant on \(\operatorname{Herm}_2(\mathbb C)\), the
positive-semidefinite Gram cone inside that vector space, and a separately
chosen Euclidean metric on the positive-definite cone.

\subsection{Conventions and the Gram morphism}

Put
\begin{equation}
 \mathscr V=\operatorname{Mat}_{3\times2}(\mathbb C),
 \qquad
 \Gamma:\mathscr V\longrightarrow\operatorname{Herm}_2(\mathbb C),
 \qquad
 \Gamma(A)=A^\dagger A.
 \label{eq:hrgc-gram-map}
\end{equation}
Write \(A=(a\ b)\), with \(a,b\in\mathbb C^3\), and let
\(\mathrm{SU}(3)\) act on the left:
\[
 U\cdot A=UA.
\]
For \(r=0,1,2\), let
\[
 \mathscr V_{(r)}=\{A\in\mathscr V:\operatorname{rank}_{\mathbb C}A=r\}.
\]
The Hermitian products are linear in the second argument.  We orient
\(\mathscr V\) by the complex orientation, in the real coordinate order
\begin{equation}
 \begin{split}
 (&\operatorname{Re}A_{11},\operatorname{Im}A_{11},
   \operatorname{Re}A_{12},\operatorname{Im}A_{12},
   \operatorname{Re}A_{21},\operatorname{Im}A_{21},\\
  &\operatorname{Re}A_{22},\operatorname{Im}A_{22},
   \operatorname{Re}A_{31},\operatorname{Im}A_{31},
   \operatorname{Re}A_{32},\operatorname{Im}A_{32}).
 \end{split}
 \label{eq:hrgc-v-orientation}
\end{equation}
On \(\operatorname{Herm}_2(\mathbb C)\), use the oriented real-linear
coordinates
\begin{equation}
 \Phi(t,x,y,z)=
 \begin{pmatrix}
  t+z&x-iy\\
  x+iy&t-z
 \end{pmatrix},
 \qquad
 \operatorname{or}_{\operatorname{Herm}}=dt\wedge dx\wedge dy\wedge dz.
 \label{eq:hrgc-pauli-coordinates}
\end{equation}
Thus, if \(\Gamma(A)_{12}=\langle a,b\rangle\), then
\begin{equation}
 \begin{aligned}
 t&=\frac12(\lVert a\rVert^2+\lVert b\rVert^2),&
 z&=\frac12(\lVert a\rVert^2-\lVert b\rVert^2),\\
 x&=\operatorname{Re}\langle a,b\rangle,&
 y&=-\operatorname{Im}\langle a,b\rangle.
 \end{aligned}
 \label{eq:hrgc-column-pauli-coordinates}
\end{equation}

\subsection{Exact orbit classification}

\begin{theorem}[The Gram quotient]
\label{thm:hrgc-gram-quotient}
The morphism \(\Gamma\) is \(\mathrm{SU}(3)\)-invariant, has image exactly
\(\operatorname{Herm}_2(\mathbb C)_{\geq0}\), and separates
\(\mathrm{SU}(3)\)-orbits.  It induces a proper homeomorphism
\begin{equation}
 \overline\Gamma:
 \mathscr V/\mathrm{SU}(3)
 \xrightarrow{\ \cong\ }
 \operatorname{Herm}_2(\mathbb C)_{\geq0}.
 \label{eq:hrgc-orbit-homeomorphism}
\end{equation}
\end{theorem}

\begin{proof}
For \(U\in\mathrm{SU}(3)\),
\[
 \Gamma(UA)=A^\dagger U^\dagger UA=A^\dagger A,
\]
so \(\Gamma\) is invariant.  Every matrix \(A^\dagger A\) is positive
semidefinite.  Conversely, for
\(H\in\operatorname{Herm}_2(\mathbb C)_{\geq0}\), let \(H^{1/2}\) be its
unique positive-semidefinite square root and define
\begin{equation}
 s(H)=
 \begin{pmatrix}
  H^{1/2}\\
  0\quad0
 \end{pmatrix}\in\operatorname{Mat}_{3\times2}(\mathbb C).
 \label{eq:hrgc-square-root-section}
\end{equation}
Then \(s(H)^\dagger s(H)=H\), proving surjectivity and giving a continuous
section.

It remains to prove exact orbit separation.  Suppose
\(A^\dagger A=B^\dagger B\).  Define
\[
 T:\operatorname{im}A\longrightarrow\operatorname{im}B,
 \qquad T(Au)=Bu.
\]
If \(Au=0\), then
\[
 0=\lVert Au\rVert^2
  =u^\dagger A^\dagger Au
  =u^\dagger B^\dagger Bu
  =\lVert Bu\rVert^2,
\]
so \(Bu=0\), and \(T\) is well defined.  For \(u,v\in\mathbb C^2\),
\[
 \langle T(Au),T(Av)\rangle
 =u^\dagger B^\dagger Bv
 =u^\dagger A^\dagger Av
 =\langle Au,Av\rangle.
\]
Thus \(T\) is a unitary isomorphism of the image subspaces.  Extend it to
some \(U_0\in\mathrm U(3)\) satisfying \(U_0A=B\).  Since
\(\operatorname{rank}B\leq2\), the orthogonal complement of
\(\operatorname{im}B\) has positive complex dimension.  Choose a unitary
map \(V\) which fixes \(\operatorname{im}B\) pointwise and whose determinant
is \((\det U_0)^{-1}\).  Then
\[
 U=VU_0\in\mathrm{SU}(3),
 \qquad UA=VB=B.
\]
Hence equal Gram matrices are exactly equal \(\mathrm{SU}(3)\)-orbits.

The action is proper because \(\mathrm{SU}(3)\) is compact; hence the orbit
space is Hausdorff.  The continuous map
\(H\mapsto[s(H)]\) is inverse to the continuous induced map
\(\overline\Gamma\), so \eqref{eq:hrgc-orbit-homeomorphism} is a
homeomorphism.  Finally,
\begin{equation}
 \lVert A\rVert_F^2=\operatorname{tr}(A^\dagger A)
 =\operatorname{tr}\Gamma(A).
 \label{eq:hrgc-properness-identity}
\end{equation}
If \(K\) is compact in the positive-semidefinite cone, then
\(\Gamma^{-1}(K)\) is closed and bounded by
\eqref{eq:hrgc-properness-identity}, hence compact.  Thus \(\Gamma\) is
proper.  For compact \(K\),
\(\overline\Gamma^{-1}(K)\) is the continuous image in the orbit space of
the compact set \(\Gamma^{-1}(K)\), so \(\overline\Gamma\) is proper as
well.
\end{proof}

\begin{proposition}[Stabilizers and orbit-type strata]
\label{prop:hrgc-stabilizers}
For \(A\in\mathscr V_{(r)}\),
\begin{equation}
 \operatorname{Stab}_{\mathrm{SU}(3)}(A)\cong\mathrm{SU}(3-r).
 \label{eq:hrgc-stabilizer-formula}
\end{equation}
The exact dimensions are
\begin{equation}
 \begin{array}{c|c|c|c|c}
 r&\dim_{\mathbb R}\mathscr V_{(r)}&\operatorname{Stab}(A)&
 \dim_{\mathbb R}(\mathrm{SU}(3)\cdot A)&
 \dim_{\mathbb R}(\mathscr V_{(r)}/\mathrm{SU}(3))\\ \hline
 2&12&\{1\}&8&4\\
 1&8&\mathrm{SU}(2)&5&3\\
 0&0&\mathrm{SU}(3)&0&0.
 \end{array}
 \label{eq:hrgc-stratum-table}
\end{equation}
Consequently the quotient has topological dimension four, but the action is
free only on the rank-two stratum.
\end{proposition}

\begin{proof}
The equality \(UA=A\) is equivalent to the assertion that \(U\) fixes
\(\operatorname{im}A\) pointwise.  Restriction to its orthogonal complement
identifies the stabilizer with \(\mathrm{SU}(3-r)\), proving
\eqref{eq:hrgc-stabilizer-formula}.  The rank-two locus is open in the
twelve-real-dimensional vector space.  The complex rank-one locus in the
space of \(3\times2\) matrices has complex dimension
\[
 1(3+2-1)=4,
\]
hence real dimension eight away from the origin.  Subtracting
\(8-\dim_{\mathbb R}\operatorname{Stab}(A)\) gives the displayed orbit and
quotient dimensions.  They agree with the target strata: positive-definite
Hermitian matrices form a four-dimensional open cone, while every nonzero
rank-one positive-semidefinite matrix has a unique expression
\(\lambda vv^\dagger\) with \(\lambda>0\) and \(v\in S^3\), modulo
\(v\sim e^{i\theta}v\).  Hence that target stratum is
\[
 (0,\infty)\times\mathbb{CP}^1
 \cong(0,\infty)\times S^2,
\]
of real dimension three.  The rank-zero stratum is the origin.
\end{proof}

\subsection{The free locus, its explicit trivialization, and orientation}

\begin{theorem}[Determinant-one completion]
\label{thm:hrgc-free-trivialization}
There is an \(\mathrm{SU}(3)\)-equivariant diffeomorphism
\begin{equation}
 \mathscr V_{(2)}
 \xrightarrow{\ \cong\ }
 \mathrm{SU}(3)\times\operatorname{Herm}_2(\mathbb C)_{>0}
 \label{eq:hrgc-free-product}
\end{equation}
under which \(\Gamma\) is projection to the second factor and the group acts
by left multiplication on the first.  Thus the principal bundle over the
free quotient has a global section and is trivial.
\end{theorem}

\begin{proof}
For \(A\in\mathscr V_{(2)}\), set
\[
 H=A^\dagger A>0,
 \qquad Q=AH^{-1/2}=(q_1\ q_2).
\]
Then \(Q^\dagger Q=I_2\), so \((q_1,q_2)\) is an ordered orthonormal
two-frame in \(\mathbb C^3\).  With \(\varepsilon_{123}=1\), define
\begin{equation}
 (q_3)_i=\sum_{j,k=1}^3\varepsilon_{ijk}
 \overline{(q_1)_j}\,\overline{(q_2)_k}.
 \label{eq:hrgc-cross-product}
\end{equation}
Antisymmetry gives
\(q_1^\dagger q_3=q_2^\dagger q_3=0\).  The Gram--Pl\"ucker identity, proved
independently in \cref{prop:hrgc-pluecker}, gives \(\lVert q_3\rVert=1\).
Finally, contraction of the two alternating tensors gives
\begin{align*}
 \det[q_1\ q_2\ q_3]
 &=\sum_{i,j,k}\varepsilon_{ijk}(q_1)_i(q_2)_j(q_3)_k\\
 &=\lVert q_1\rVert^2\lVert q_2\rVert^2
   -|q_1^\dagger q_2|^2=1.
\end{align*}
Thus \(U_A=[q_1\ q_2\ q_3]\in\mathrm{SU}(3)\), and
\[
 A=U_A s(H),
\]
where \(s\) is \eqref{eq:hrgc-square-root-section}.  The positive square
root, inverse square root, and the polynomial-conjugate formula
\eqref{eq:hrgc-cross-product} are smooth on the positive-definite locus.
Consequently
\[
 A\longmapsto(U_A,A^\dagger A),
 \qquad
 (U,H)\longmapsto Us(H)
\]
are smooth inverse maps.  Their equivariance is immediate.
\end{proof}

The order of factors in \eqref{eq:hrgc-free-product} does not by itself fix
an orientation.  We now fix and calculate it.  Give
\(\mathfrak{su}(3)\) the ordered real basis
\begin{equation}
 \begin{gathered}
 X_{12},Y_{12},X_{13},Y_{13},X_{23},Y_{23},D_1,D_2,\\
 X_{jk}=E_{jk}-E_{kj},\qquad
 Y_{jk}=i(E_{jk}+E_{kj}),\\
 D_1=i\operatorname{diag}(1,-1,0),\qquad
 D_2=i\operatorname{diag}(0,1,-1).
 \end{gathered}
 \label{eq:hrgc-su3-oriented-basis}
\end{equation}
At
\[
 A_0=\begin{pmatrix}1&0\\0&1\\0&0\end{pmatrix},
\]
use the oriented Hermitian basis
\begin{equation}
 K_t=I_2,\quad
 K_x=\begin{pmatrix}0&1\\1&0\end{pmatrix},\quad
 K_y=\begin{pmatrix}0&-i\\i&0\end{pmatrix},\quad
 K_z=\begin{pmatrix}1&0\\0&-1\end{pmatrix}
 \label{eq:hrgc-hermitian-oriented-basis}
\end{equation}
and horizontal lifts \(B_\alpha=\tfrac12A_0K_\alpha\).  Directly,
\(d\Gamma_{A_0}(B_\alpha)=K_\alpha\).  Relative to
\eqref{eq:hrgc-v-orientation}, the twelve columns
\[
 (X_{12}A_0,Y_{12}A_0,X_{13}A_0,Y_{13}A_0,
 X_{23}A_0,Y_{23}A_0,D_1A_0,D_2A_0,B_t,B_x,B_y,B_z)
\]
form the exact real coordinate matrix
\begingroup
\scriptsize
\setlength{\arraycolsep}{2.4pt}
\begin{equation}
 M_{\mathrm{or}}:=\left(
 \begin{array}{rrrrrrrrrrrr}
0&0&0&0&0&0&0&0&\frac12&0&0&\frac12\\
0&0&0&0&0&0&1&0&0&0&0&0\\
1&0&0&0&0&0&0&0&0&\frac12&0&0\\
0&1&0&0&0&0&0&0&0&0&-\frac12&0\\
-1&0&0&0&0&0&0&0&0&\frac12&0&0\\
0&1&0&0&0&0&0&0&0&0&\frac12&0\\
0&0&0&0&0&0&0&0&\frac12&0&0&-\frac12\\
0&0&0&0&0&0&-1&1&0&0&0&0\\
0&0&-1&0&0&0&0&0&0&0&0&0\\
0&0&0&1&0&0&0&0&0&0&0&0\\
0&0&0&0&-1&0&0&0&0&0&0&0\\
0&0&0&0&0&1&0&0&0&0&0&0
 \end{array}
 \right).
 \label{eq:hrgc-orientation-matrix}
\end{equation}
\endgroup
Laplace expansion gives
\begin{equation}
 \det M_{\mathrm{or}}=\frac12>0.
 \label{eq:hrgc-orientation-determinant}
\end{equation}
Since both factors in \eqref{eq:hrgc-free-product} are connected, the sign is
constant on the free locus.  Therefore the quotient orientation determined
by
\[
 \operatorname{or}_{\mathscr V}
 =\operatorname{or}_{\mathrm{SU}(3)}\wedge
  \operatorname{or}_{\mathscr V_{(2)}/\mathrm{SU}(3)}
\]
agrees under \(\Gamma\) with
\(dt\wedge dx\wedge dy\wedge dz\).

\subsection{Lorentz determinant and Pl\"ucker norm}

For \(H=\Phi(t,x,y,z)\), direct expansion gives
\begin{equation}
 \operatorname{tr}H=2t,
 \qquad
 \det H=t^2-x^2-y^2-z^2,
 \label{eq:hrgc-lorentz-determinant}
\end{equation}
and the eigenvalues are
\(t\pm\sqrt{x^2+y^2+z^2}\).  Hence
\begin{equation}
 \begin{aligned}
 H>0&\Longleftrightarrow t>\sqrt{x^2+y^2+z^2},\\
 H\geq0&\Longleftrightarrow t\geq\sqrt{x^2+y^2+z^2},\\
 \operatorname{rank}H=1&\Longleftrightarrow
 t=\sqrt{x^2+y^2+z^2}>0.
 \end{aligned}
 \label{eq:hrgc-positive-cone-inequalities}
\end{equation}
Thus \(\det\) has signature \((1,3)\) when positive directions are listed
first, while \(-\det\) has signature \((3,1)\).  The quotient image is not
all of this Lorentz vector space: it is precisely the closed future causal
cone, and its free locus is the open future timelike cone.

\begin{proposition}[Gram--Pl\"ucker identity]
\label{prop:hrgc-pluecker}
For \(A=(a\ b)\),
\begin{equation}
 \det\Gamma(A)
 =\lVert a\rVert^2\lVert b\rVert^2-|\langle a,b\rangle|^2
 =\sum_{1\leq i<j\leq3}|a_ib_j-a_jb_i|^2
 =\lVert a\wedge b\rVert^2.
 \label{eq:hrgc-pluecker-identity}
\end{equation}
\end{proposition}

\begin{proof}
Expanding gives
\begin{align*}
 \lVert a\rVert^2\lVert b\rVert^2-|\langle a,b\rangle|^2
 &=\sum_{i,j}|a_i|^2|b_j|^2
   -\sum_{i,j}\overline{a_i}b_i a_j\overline{b_j}\\
 &=\sum_{i<j}
 \bigl(|a_i|^2|b_j|^2+|a_j|^2|b_i|^2
 -\overline{a_i}b_i a_j\overline{b_j}
 -\overline{a_j}b_j a_i\overline{b_i}\bigr)\\
 &=\sum_{i<j}|a_ib_j-a_jb_i|^2.
\end{align*}
The last expression is the squared Hermitian norm of \(a\wedge b\) in the
ordered basis
\((e_1\wedge e_2,e_1\wedge e_3,e_2\wedge e_3)\).
\end{proof}

The determinant is only one invariant and does not classify orbits.  For
example,
\[
 A_1=\begin{pmatrix}1&0\\0&1\\0&0\end{pmatrix},
 \qquad
 A_2=\begin{pmatrix}\sqrt2&0\\0&1/\sqrt2\\0&0\end{pmatrix}
\]
satisfy \(\det\Gamma(A_1)=\det\Gamma(A_2)=1\), but their Gram matrices are
\(I_2\) and \(\operatorname{diag}(2,1/2)\), so
\cref{thm:hrgc-gram-quotient} shows that they lie in distinct orbits.

\subsection{The four-coefficient form is additional input}

Let \(r=(r_1,r_2,r_3,r_4)^{\mathsf T}\in\mathbb R^4\), and define the
orthogonal Hadamard coordinates
\begin{equation}
 \begin{pmatrix}h\\x\\y\\z\end{pmatrix}
 =\frac12
 \begin{pmatrix}
 1&1&1&1\\
 1&1&-1&-1\\
 1&-1&1&-1\\
 1&-1&-1&1
 \end{pmatrix}
 \begin{pmatrix}r_1\\r_2\\r_3\\r_4\end{pmatrix}.
 \label{eq:hrgc-hadamard-transform}
\end{equation}
The proposed coefficient-space form
\begin{equation}
 Q_4(r)=\frac12\left(\sum_{i=1}^4r_i\right)^2-
 \sum_{i=1}^4r_i^2
 \label{eq:hrgc-q4-definition}
\end{equation}
satisfies the exact identity
\begin{equation}
 Q_4(r)=h^2-x^2-y^2-z^2.
 \label{eq:hrgc-q4-lorentz}
\end{equation}
Indeed, orthogonality gives
\(\sum r_i^2=h^2+x^2+y^2+z^2\) and
\(\sum r_i=2h\).

The same Hermitian matrix can be written in the alternative convention
\begin{equation}
 \Psi(h,x,y,z)=
 \begin{pmatrix}h+y&x+iz\\x-iz&h-y\end{pmatrix}.
 \label{eq:hrgc-pro-hermitian-convention}
\end{equation}
It agrees exactly with \(\Phi\) under
\[
 (t,x,y_{\Phi},z_{\Phi})=(h,x,-z,y).
\]
The transformation \((y,z)\mapsto(-z,y)\) has determinant \(+1\), so the
two conventions give the same matrix, Lorentz form, and orientation.

Nothing in the \(\mathrm{SU}(3)\) orbit calculation selects the coefficient
\(1/2\).  More generally,
\begin{equation}
 Q_{4,\lambda}(r)=\lambda\left(\sum_ir_i\right)^2-\sum_ir_i^2
 =(4\lambda-1)h^2-x^2-y^2-z^2.
 \label{eq:hrgc-q4-lambda-family}
\end{equation}
Every \(\lambda>1/4\) therefore has Lorentz signature \((1,3)\), after a
rescaling of \(h\).  The value \(\lambda=1/2\) is the unique member for which
the fixed orthogonal coordinates in \eqref{eq:hrgc-hadamard-transform} have
equal unit coefficients.  This is a normalization choice in coefficient
space, not a theorem forced by the Gram quotient.

Nor does a nonnegative coefficient tuple automatically determine a point of
the positive Gram cone.  The tuple \((1,0,0,0)\) has
\(Q_4=-1/2\).  A coefficient tuple enters the Gram quotient through
\eqref{eq:hrgc-pro-hermitian-convention} only when that Hermitian matrix is
positive semidefinite, equivalently when \(h\geq0\) and \(Q_4\geq0\).
The hypersurface \(z=0\) is a real-symmetric coefficient slice, but that
fact alone does not supply a coordinatewise nonnegative embedding of a
three-component mass domain.

\subsection{The compact representation quotient}

\begin{theorem}[The \(S^{12}\) representation quotient]
\label{thm:hrgc-s12-ball-quotient}
Let
\[
 S^{12}_{\mathrm{col}}=\mathscr V^+=\mathscr V\cup\{\infty\}
\]
be the one-point compactification, and extend the \(\mathrm{SU}(3)\) action
by fixing \(\infty\).  Then
\begin{equation}
 S^{12}_{\mathrm{col}}/\mathrm{SU}(3)
 \cong
 \bigl(\operatorname{Herm}_2(\mathbb C)_{\geq0}\bigr)^+
 \cong B^4.
 \label{eq:hrgc-s12-ball}
\end{equation}
The orbit-type strata are
\begin{equation}
 \begin{array}{c|c|c}
 \text{source stratum}&\text{quotient stratum}&\text{stabilizer}\\ \hline
 \operatorname{rank}A=2&\operatorname{int}B^4&\{1\}\\
 \operatorname{rank}A=1&\partial B^4\setminus\{p_0,p_\infty\}&
   \mathrm{SU}(2)\\
 A=0&\{p_0\}&\mathrm{SU}(3)\\
 A=\infty&\{p_\infty\}&\mathrm{SU}(3).
 \end{array}
 \label{eq:hrgc-compact-strata}
\end{equation}
The ball identification is topological on the full stratified quotient.  On
its smooth interior, it may be chosen to carry the quotient orientation to
the orientation fixed in \eqref{eq:hrgc-pauli-coordinates}.
\end{theorem}

\begin{proof}
As a real vector space, \(\mathscr V\cong\mathbb R^{12}\), so its one-point
compactification is \(S^{12}\).  On
\(\mathscr V\cong\mathbb C^3\oplus\mathbb C^3\), the action of
\(U\in\mathrm{SU}(3)\) has complex determinant \((\det U)^2=1\), hence
real determinant one; it preserves the standard orientation before and after
one-point compactification.

The properness proved from \eqref{eq:hrgc-properness-identity} implies that
\(\Gamma\) extends continuously to
\[
 \Gamma^+:\mathscr V^+\longrightarrow
 \bigl(\operatorname{Herm}_2(\mathbb C)_{\geq0}\bigr)^+,
 \qquad \Gamma^+(\infty)=\infty.
\]
It separates the finite orbits by \cref{thm:hrgc-gram-quotient}, and the
added point is a distinct fixed orbit.  Hence it induces a continuous
bijection
\[
 \mathscr V^+/\mathrm{SU}(3)
 \longrightarrow
 \bigl(\operatorname{Herm}_2(\mathbb C)_{\geq0}\bigr)^+.
\]
The source is compact Hausdorff because a compact group acts on the compact
Hausdorff space \(\mathscr V^+\).  The target is Hausdorff because the
positive-semidefinite cone is locally compact Hausdorff.  The continuous
bijection is therefore a homeomorphism.

By \eqref{eq:hrgc-positive-cone-inequalities}, the map
\[
 (t,\mathbf x)\longmapsto(t-\lVert\mathbf x\rVert,\mathbf x)
\]
is a homeomorphism from the closed future cone to the closed half-space
\(\mathbb R_{\geq0}\times\mathbb R^3\).  For \(u=(s,\mathbf x)\) in that half-space,
inverse stereographic projection is
\begin{equation}
 u\longmapsto
 \frac{1}{1+\lVert u\rVert^2}
 \bigl(2s,2\mathbf x,\lVert u\rVert^2-1\bigr)
 \in S^4\subset\mathbb R^5.
 \label{eq:hrgc-halfspace-stereographic}
\end{equation}
Its image is the closed hemisphere with first coordinate nonnegative, minus
the north pole.  The point at infinity supplies that pole.  A closed
hemisphere is a closed four-ball, proving \eqref{eq:hrgc-s12-ball}.

Positive-definite matrices correspond to the open half-space and hence the
ball interior.  Nonzero rank-one matrices correspond to
\(\{0\}\times(\mathbb R^3\setminus\{0\})\).  The zero matrix and the point at
infinity give two distinct endpoints in the compactification of that
three-dimensional boundary.  Combining this with
\cref{prop:hrgc-stabilizers} proves \eqref{eq:hrgc-compact-strata}.  The
orientation statement on the connected interior is the calculation following
\eqref{eq:hrgc-orientation-determinant}.
\end{proof}

The notation \(S^{12}_{\mathrm{col}}\) records the construction.  There is a
canonical pointed homeomorphism of representation spheres
\[
 (\mathbb C^3)^+\wedge(\mathbb C^3)^+
 \cong(\mathbb C^3\oplus\mathbb C^3)^+
 =S^{12}_{\mathrm{col}}.
\]
If a different based six-sphere \(X\) is used, a based homeomorphism
\(X\cong(\mathbb C^3)^+\) transports the action to \(X\wedge X\).  The
transported action depends on that choice.  Without an independently proved
equivariance statement, it is not induced by a torus fibration, cusp
monodromy, period matrix, or fibrewise flux on \(X\).

\subsection{Exact gauge-theoretic boundary}

\begin{theorem}[No native nonzero \(S^4\) sector from the Gram quotient]
\label{thm:hrgc-no-native-s4-sector}
The construction above supplies no native principal \(\mathrm{SU}(3)\)
bundle of nonzero second Chern number on \(S^4\).  More precisely:
\begin{enumerate}
\item the orbit projection over the rank-two locus is the product bundle in
\eqref{eq:hrgc-free-product}, so its second Chern class is zero;
\item on the rank-one boundary the orbit projection is not a principal
\(\mathrm{SU}(3)\)-bundle, because its stabilizer is \(\mathrm{SU}(2)\);
\item the full compact quotient is \(B^4\), not \(S^4\); and
\item replacing \(B^4\) by the one-point compactification of its interior is
a different construction whose clutching class is additional data.
\end{enumerate}
\end{theorem}

\begin{proof}
The first two assertions are
\cref{thm:hrgc-free-trivialization,prop:hrgc-stabilizers}; the third is
\cref{thm:hrgc-s12-ball-quotient}.  For the fourth, the matrix logarithm is a
diffeomorphism
\[
 \log:\operatorname{Herm}_2(\mathbb C)_{>0}
 \xrightarrow{\ \cong\ }
 \operatorname{Herm}_2(\mathbb C)\cong\mathbb R^4.
\]
Its one-point compactification is therefore \(S^4\), but this operation
collapses the entire finite rank-one boundary of the full Gram quotient into
the new point at infinity.  Principal \(\mathrm{SU}(3)\)-bundles on this
new \(S^4\) are obtained by clutching two trivial bundles over hemispheres
with a map
\[
 \kappa:S^3\longrightarrow\mathrm{SU}(3).
\]
The fibration
\(\mathrm{SU}(2)\to\mathrm{SU}(3)\to S^5\), together with
\(\pi_4(S^5)=\pi_3(S^5)=0\), shows that
\[
 \pi_3(\mathrm{SU}(3))\cong\pi_3(\mathrm{SU}(2))
 =\pi_3(S^3)\cong\mathbb Z.
\]
Choose the generator so that its associated bundle has second Chern number
\(+1\).  A clutching class \(k\in\mathbb Z\) then has second Chern number
\(k\).  Extending the explicit product framing in
\cref{thm:hrgc-free-trivialization} by the same framing gives the constant
clutching map and \(k=0\).  Neither \(A^\dagger A\) nor
\(\lVert a\wedge b\rVert^2\) selects a nonconstant \(\kappa\).  Thus
\(k\ne0\) is valid additional gauge data, but it is not produced by the Gram
construction.
\end{proof}

There are two further cohomological checks.  If a candidate total space is
diffeomorphic to \(S^6\), then
\[
 H^2(S^6;\mathbb Z)=H^4(S^6;\mathbb Z)=0.
\]
Consequently no nonzero fibrewise degree-two class can be the restriction of
a degree-two class on that total space, and no \(c_2\) on the total space can
restrict to a nonzero degree-four class on a fibre.  This does not forbid
independently chosen higher characteristic data; it proves exactly that the
degree-two and degree-four classes in question do not extend.

Likewise, for every continuous map \(f:S^{12}\to S^4\) and generator
\(u_4\in H^4(S^4;\mathbb Z)\),
\[
 f^*u_4=0
\]
because \(H^4(S^{12};\mathbb Z)=0\).  Therefore the pullback along any such
map of an \(\mathrm{SU}(3)\) bundle with nonzero \(c_2=ku_4\) has
\(c_2=0\).  This statement does not assert that every intrinsic
\(\mathrm{SU}(3)\) bundle on \(S^{12}\) is trivial; it says that an
\(S^4\) second-Chern charge is not inherited by pullback to the
representation sphere.

\subsection{The log-Euclidean metric and Yang--Mills theory are extra choices}

The logarithm identifies the open positive cone with a real vector space.
One may therefore declare
\begin{equation}
 X=\log H,
 \qquad
 g_{\log}=\frac12\operatorname{tr}(dX\,dX),
 \label{eq:hrgc-log-metric}
\end{equation}
which is a complete flat Euclidean metric in the \(X\)-coordinates.  This
metric is not induced by the Frobenius metric on \(\mathscr V\).  To see the
difference without discarding any scale factor, set
\[
 H(s)=sI_2,
 \qquad
 A(s)=\sqrt{s}\begin{pmatrix}1&0\\0&1\\0&0\end{pmatrix},
 \qquad s>0.
\]
For \(\xi\in\mathfrak{su}(3)\),
\[
 \operatorname{Re}\operatorname{tr}
 \bigl(\dot A^\dagger\xi A\bigr)
 =\frac{\dot s}{2}
 \operatorname{Re}(\xi_{11}+\xi_{22})=0,
\]
so \(\dot A\) is horizontal for the Frobenius quotient.  Hence
\begin{equation}
 \lVert\dot H\rVert_{\mathrm{quot}}^2
 =\lVert\dot A\rVert_F^2=\frac{\dot s^2}{2s}.
 \label{eq:hrgc-frobenius-ray}
\end{equation}
By contrast, \(X(s)=(\log s)I_2\), and
\begin{equation}
 \lVert\dot H\rVert_{\log}^2
 =\frac12\operatorname{tr}(\dot X^2)
 =\frac{\dot s^2}{s^2}.
 \label{eq:hrgc-log-ray}
\end{equation}
The ratio of \eqref{eq:hrgc-frobenius-ray} and
\eqref{eq:hrgc-log-ray} depends on \(s\), so no global constant rescaling
identifies the metrics.  The Lorentz form
\eqref{eq:hrgc-lorentz-determinant} is a third object and is indefinite.

After choosing \eqref{eq:hrgc-log-metric}, one may define classical
Yang--Mills theory on the trivial product bundle supplied by
\cref{thm:hrgc-free-trivialization}.  The additional choice of the zero
connection has \(F_0=0\), so it solves the nonlinear classical equation
\(d_0^*F_0=0\); its Coulomb Hessian is the transverse free Laplacian.  These
are the standard zero-charge vacuum statements on a chosen Euclidean
\(\mathbb R^4\).  The Gram morphism itself supplies no nonzero clutching
class, no nonflat connection, no interacting Hamiltonian, no quantum
Hilbert-space identification, and no mass-gap estimate.  Thus the exact
logical implication is
\[
 \boxed{
 \begin{array}{c}
 \text{four-dimensional Gram/Lorentz quotient: proved;}\\
 S^{12}_{\mathrm{col}}/\mathrm{SU}(3)\cong B^4: \text{ proved;}\\
 \text{nonzero native }\mathrm{SU}(3)\text{ sector on }S^4:
 \text{ not produced;}\\
 \text{interacting or quantum Yang--Mills inheritance: not produced.}
 \end{array}}
\]


\section{The marked rank-six toroidal lattice and the Weyl carrier}
\label{sec:rank-six-toroidal-lattice}

The rank-six lattice in this section is the first homology of the marked
rank-three torus.  It is not the rank-six degree-four defect group introduced
in \cref{subsec:distinct-defect-lattice} below.  Put
\[
 L_3=\mathbb Z^3,\qquad
 \bar L_3=\mathbb Z^3,\qquad
 \Lambda_3=L_3\oplus\bar L_3=H_1(T^6;\mathbb Z),
 \qquad u=(1,1,1)^{\mathsf T}.
\]
For \(\sigma\in S_3\), let \(R_\sigma e_i=e_{\sigma(i)}\).  The retained
Weyl action on \(\Lambda_3\) is the diagonal action
\begin{equation}
 \sigma(a,b)=(R_\sigma a,R_\sigma b).                    \label{eq:rank6-s3-action}
\end{equation}
The monodromy, nilpotent part, filling map, and integral symplectic form are
\begin{align}
 M_3(a,b)&=(a+b,b),& N_3&=M_3-I, & c_3(a,b)&=b,
                                                        \label{eq:rank6-monodromy}\\
 \Omega_3((a,b),(a',b'))&=a^{\mathsf T}b'-b^{\mathsf T}a'.
                                                        \label{eq:rank6-symplectic}
\end{align}

\begin{proposition}[Integral Weyl and monodromy structure]
\label{prop:rank6-weyl-monodromy}
The \(S_3\)-action commutes with \(M_3\), preserves \(c_3\) equivariantly,
and preserves \(\Omega_3\).  Moreover
\[
 N_3^2=0,\qquad
 \ker N_3=\operatorname{im}N_3=L_3\oplus0.
\]
As an integral \(S_3\)-module,
\[
 \Lambda_3\cong\mathbb Z^3_{\rm perm}\oplus
                    \mathbb Z^3_{\rm perm};
\]
over \(\mathbb Q\) this is
\begin{equation}
 \Lambda_3\otimes\mathbb Q
 \cong2\,\mathbf1\oplus2\,{\rm Std}.                  \label{eq:rank6-rational-decomposition}
\end{equation}
Its character on the conjugacy classes represented by
\(1,(12),(123)\) is respectively \((6,2,0)\).
\end{proposition}

\begin{proof}
Every permutation matrix is orthogonal and fixes \(u\).  Substitution in
\eqref{eq:rank6-monodromy} shows that the diagonal action commutes with the
shear and that \(c_3(\sigma(a,b))=R_\sigma c_3(a,b)\).  Orthogonality gives
\[
 \Omega_3(\sigma(a,b),\sigma(a',b'))
 =a^{\mathsf T}R_\sigma^{\mathsf T}R_\sigma b'
  -b^{\mathsf T}R_\sigma^{\mathsf T}R_\sigma a'
 =\Omega_3((a,b),(a',b')).
\]
Also \(N_3(a,b)=(b,0)\), proving the two nilpotent assertions.  Each of the
two displayed bases is a three-point permutation orbit.  A permutation
module \(\mathbb Q^3\) is \(\mathbf1\oplus{\rm Std}\); doubling gives
\eqref{eq:rank6-rational-decomposition}.  The traces are twice the numbers
of fixed points of a permutation acting on three letters, hence \(6,2,0\).
\end{proof}

\subsection{The exact quotient and its nonsplitting as a Weyl lattice}

Let \(A_F=B_F=\mathbb Z^2\), and retain the matrices
\begin{equation}
 B_0=\begin{pmatrix}0&1\\-1&0\end{pmatrix},\qquad
 P=\begin{pmatrix}1&-1&0\\0&1&-1\end{pmatrix},\qquad
 Q=B_0^{-1}P=\begin{pmatrix}0&-1&1\\1&-1&0\end{pmatrix}.
                                                        \label{eq:rank6-pq}
\end{equation}
Define
\begin{equation}
 \Phi=P\oplus Q:\Lambda_3\longrightarrow A_F\oplus B_F,
 \qquad \Phi(a,b)=(Pa,Qb),                              \label{eq:rank6-phi}
\end{equation}
and
\[
 M_F(x,y)=(x+B_0y,y),\qquad c_F(x,y)=y.
\]

\begin{theorem}[Kernel, weight lattices, and integral nonsplitting]
\label{thm:rank6-exact-sequence}
The map \(\Phi\) is surjective and fits into an exact sequence of integral
\(S_3\)-modules
\begin{equation}
 0\longrightarrow K\longrightarrow\Lambda_3
 \xrightarrow{\ \Phi\ }A_F\oplus B_F\longrightarrow0,
 \qquad
 K=\mathbb Z(u,0)\oplus\mathbb Z(0,u).                 \label{eq:rank6-exact}
\end{equation}
The action on \(A_F\) and \(B_F\) is the uniquely induced action satisfying
\begin{equation}
 \rho_A(\sigma)P=PR_\sigma,qquad
 \rho_B(\sigma)Q=QR_\sigma.                            \label{eq:rank6-target-action}
\end{equation}
It obeys
\[
 \rho_A(\sigma)B_0=B_0\rho_B(\sigma),qquad
 \Phi M_3=M_F\Phi,qquad c_F\Phi=Qc_3.
\]
Both \(A_F\) and \(B_F\), with these actions, are the \(A_2\) weight
lattice \(P(A_2)=\mathbb Z^3/\mathbb Zu\).  They are not the root lattice
\[
 Q(A_2)=\{a\in\mathbb Z^3:u^{\mathsf T}a=0\},
\]
whose image has index three in either quotient.  Sequence
\eqref{eq:rank6-exact} splits in the category of abelian groups, but it has
no integral \(S_3\)-equivariant splitting.
\end{theorem}

\begin{proof}
The two-by-two minors of \(P\) include \(1\), so \(P\) is surjective; the
identity \(Q=B_0^{-1}P\), with \(B_0\in\mathrm{SL}_2(\mathbb Z)\), gives
the same conclusion for \(Q\).  Direct solution of \(Pa=0\) gives
\(a_1=a_2=a_3\), and therefore
\(\ker P=\ker Q=\mathbb Zu\).  This proves
\eqref{eq:rank6-exact}, and \(K\) is fixed pointwise by \(S_3\).

Because \(P\) and \(Q\) are quotient maps with invariant kernel, the formulas
\eqref{eq:rank6-target-action} define unique integral actions.  From
\(B_0Q=P\) and the surjectivity of \(Q\),
\[
 \rho_A(\sigma)B_0Q=\rho_A(\sigma)P
 =PR_\sigma=B_0QR_\sigma=B_0\rho_B(\sigma)Q,
\]
which proves the first compatibility.  The other two are the direct
calculations
\[
 \Phi M_3(a,b)=(Pa+Pb,Qb)
 =(Pa+B_0Qb,Qb)=M_F\Phi(a,b)
\]
and \(c_F\Phi(a,b)=Qb=Qc_3(a,b)\).

The first isomorphism theorem identifies each quotient with
\(\mathbb Z^3/\mathbb Zu\).  In the root basis
\(\alpha_1=e_1-e_2,\alpha_2=e_2-e_3\),
\begin{equation}
 P(\alpha_1\ \alpha_2)
 =\begin{pmatrix}2&-1\\-1&2\end{pmatrix}.              \label{eq:rank6-cartan-index}
\end{equation}
The determinant is three, proving that the root lattice maps with index
three into the weight lattice.  An explicit abelian-group section of \(P\)
is
\[
 S=\begin{pmatrix}1&0\\0&0\\0&-1\end{pmatrix},
 \qquad PS=I_2;
\]
an explicit section of \(Q\) is \(SB_0\).  Hence
\((x,y)\mapsto(Sx,SB_0y)\) splits \(\Phi\) as abelian groups.

It remains to prove the integral equivariant obstruction.  Over \(\mathbb Q\)
the invariant line \(\mathbb Qu\) has the unique isotypic complement
\(u^\perp=\{a:u^{\mathsf T}a=0\}\); there are no equivariant maps from
\({\rm Std}\) to \(\mathbf1\).  Thus any rational \(S_3\)-equivariant
complement to \(K\otimes\mathbb Q\) in \(\Lambda_3\otimes\mathbb Q\) must
be \(u^\perp\oplus u^\perp\).  Its integral intersection is
\(Q(A_2)\oplus Q(A_2)\), whose image in \(A_F\oplus B_F\) has index
\(3^2=9\) by \eqref{eq:rank6-cartan-index}.  It is not surjective, so an
integral equivariant section cannot exist.
\end{proof}

\begin{proposition}[The symplectic position of the kernel]
\label{prop:rank6-kernel-pairing}
In the ordered basis \(((u,0),(0,u))\),
\begin{equation}
 \Omega_3|_K=\begin{pmatrix}0&3\\-3&0\end{pmatrix}.     \label{eq:rank6-k-pairing}
\end{equation}
Moreover
\[
 K^{\perp_{\Omega_3}}=Q(A_2)\oplus Q(A_2),\qquad
 [\Lambda_3:K\oplus K^{\perp_{\Omega_3}}]=9.
\]
In the root bases \((\alpha_1,\alpha_2)\) in the two summands, the restricted
form is
\begin{equation}
 \Omega_3|_{K^{\perp_{\Omega_3}}}
 =\begin{pmatrix}0&C\\-C&0\end{pmatrix},
 \qquad C=\begin{pmatrix}2&-1\\-1&2\end{pmatrix}.       \label{eq:rank6-root-pairing}
\end{equation}
Both \(K\) and its symplectic orthogonal are preserved by \(S_3\) and
\(M_3\).  In particular, \(\Omega_3\) does not descend through \(\Phi\) to
a form on \(A_F\oplus B_F\).
\end{proposition}

\begin{proof}
Equation \eqref{eq:rank6-k-pairing} follows from \(u^{\mathsf T}u=3\).
Orthogonality to \((u,0)\) and \((0,u)\) is respectively the pair of
conditions \(u^{\mathsf T}b=0\) and \(u^{\mathsf T}a=0\).  The determinant
of the three-column matrix \((u\ \alpha_1\ \alpha_2)\) has absolute value
three; taking both summands gives index nine.  Finally
\((\alpha_i^{\mathsf T}\alpha_j)_{ij}=C\), which proves
\eqref{eq:rank6-root-pairing}.  The preservation statements follow from
\cref{prop:rank6-weyl-monodromy}.  A bilinear form descending through a
quotient must annihilate the kernel in each argument, whereas
\eqref{eq:rank6-k-pairing} is nonzero.
\end{proof}

The simultaneous quotient--kernel matrices
\[
 A=\begin{pmatrix}P\\e_3^{\mathsf T}\end{pmatrix},
 \qquad B=\begin{pmatrix}Q\\e_3^{\mathsf T}\end{pmatrix},
 \qquad J=\operatorname{diag}(A,B)
\]
are unimodular and split the complex-torus lattice and its monodromy.  They
do not preserve the particular marked alternating form
\eqref{eq:rank6-symplectic}.  In fact, with
\(\Omega_3=\begin{psmallmatrix}0&I_3\\-I_3&0\end{psmallmatrix}\),
\begin{equation}
 J^{\mathsf T}\Omega_3J
 =\begin{pmatrix}0&H\\-H^{\mathsf T}&0\end{pmatrix},
 \qquad
 H=A^{\mathsf T}B=
 \begin{pmatrix}0&-1&1\\1&0&-1\\-1&1&1\end{pmatrix}\ne I_3.
                                                        \label{eq:rank6-j-not-symplectic}
\end{equation}
Thus this product splitting is an unpolarized complex-torus isomorphism; it
cannot be called an isomorphism of principally polarized tori for the
standard marked form without an additional polarization comparison.

\subsection{Six chambers, six formal rays, and their exact incidence}

Let \(E_{ii}\) be the diagonal matrix unit and put
\[
 \nu_i=E_{ii}-\frac13I_3.
\]
For a one-line permutation
\(\pi=(\pi_1\pi_2\pi_3)\in S_3\), the marked \(A_2\) chamber is
\begin{equation}
 \mathcal C_\pi
 =\mathbb R_{\geq0}\nu_{\pi_1}
  +\mathbb R_{\geq0}(-\nu_{\pi_3}).                    \label{eq:rank6-chamber}
\end{equation}
These are the images under the angular quotient \(P\) of the six Kuhn
tetrahedra incident to a marked vertex.  Notice that both rays in
\eqref{eq:rank6-chamber} are angular: the second is the image of a negative
angular vector, not of a deck vector under \(Q\).

Let
\[
 \mathscr C=\mathbb Z\langle c_\pi:\pi\in S_3\rangle
 \cong\mathbb Z[S_3]
\]
be the chamber module and let
\[
 \mathscr R=\mathbb Z\langle r_1^+,r_2^+,r_3^+,
                                  r_1^-,r_2^-,r_3^-\rangle
\]
be the formal oriented-ray module.  The action permutes indices.  Define
\begin{align}
 I:\mathscr C&\longrightarrow\mathscr R,&
 I(c_\pi)&=r_{\pi_1}^++r_{\pi_3}^-,                   \label{eq:rank6-incidence-map}\\
 d:\mathscr R&\longrightarrow\mathbb Z,&
 d(r_i^+)&=1,& d(r_i^-)&=-1.                           \label{eq:rank6-ray-augmentation}
\end{align}

\begin{theorem}[Integral chamber--ray incidence]
\label{thm:rank6-incidence}
There is an exact sequence of integral \(S_3\)-modules
\begin{equation}
 0\longrightarrow\mathbb Z_{\rm sgn}
 \xrightarrow{\ 1\mapsto\varepsilon\ }\mathscr C
 \xrightarrow{\ I\ }\mathscr R
 \xrightarrow{\ d\ }\mathbb Z_{\rm triv}
 \longrightarrow0,                                    \label{eq:rank6-incidence-exact}
\end{equation}
where
\begin{equation}
 \varepsilon=\sum_{\pi\in S_3}\operatorname{sgn}(\pi)c_\pi.
                                                        \label{eq:rank6-epsilon}
\end{equation}
Use the chamber order
\[
 (123),(132),(213),(231),(312),(321)
\]
and the ray order
\[
 (r_1^+,r_2^+,r_3^+,r_1^-,r_2^-,r_3^-).
\]
Then the matrix of \(I\) is
\begin{equation}
 [I]=
 \begin{pmatrix}
 1&1&0&0&0&0\\
 0&0&1&1&0&0\\
 0&0&0&0&1&1\\
 0&0&0&1&0&1\\
 0&1&0&0&1&0\\
 1&0&1&0&0&0
 \end{pmatrix}.                                        \label{eq:rank6-incidence-matrix}
\end{equation}
With the basis vectors orthonormal, the chamber incidence pairing is
\begin{equation}
 I^{\mathsf T}I=
 \begin{pmatrix}
 2&1&1&0&0&0\\
 1&2&0&0&1&0\\
 1&0&2&1&0&0\\
 0&0&1&2&0&1\\
 0&1&0&0&2&1\\
 0&0&0&1&1&2
 \end{pmatrix},                                        \label{eq:rank6-incidence-pairing}
\end{equation}
and its integral radical is \(\mathbb Z\varepsilon\).
\end{theorem}

\begin{proof}
The six possible ordered pairs \((\pi_1,\pi_3)\) are precisely the pairs
\((i,j)\) with \(i\ne j\), giving \eqref{eq:rank6-incidence-matrix}.
After replacing each negative-ray basis vector by its negative, this is the
oriented vertex--edge incidence map of the connected bipartite graph having
vertices \(r_i^+,r_j^-\) and one edge for every \(i\ne j\).  That graph has
six vertices and six edges and is a single six-cycle.  Its integral cellular
chain complex is exact except for its rank-one first homology.  Returning to
the original signs identifies the last vertex augmentation with \(d\), and
the cycle generator with \eqref{eq:rank6-epsilon}.  Left multiplication by
\(\sigma\) sends \(\varepsilon\) to
\(\operatorname{sgn}(\sigma)\varepsilon\), proving the module labels in
\eqref{eq:rank6-incidence-exact}.  Direct multiplication of
\eqref{eq:rank6-incidence-matrix} gives
\eqref{eq:rank6-incidence-pairing}.  Since the Euclidean ray pairing is
positive definite, \(\ker(I^{\mathsf T}I)=\ker I=\mathbb Z\varepsilon\).
\end{proof}

The chamber set is one free \(S_3\)-orbit, so its permutation character is
\((6,0,0)\).  By contrast, the six displayed basis vectors of \(\Lambda_3\)
are two three-point orbits and have character \((6,2,0)\) by
\cref{prop:rank6-weyl-monodromy}.  Thus the six chambers and the six basis
directions are not isomorphic as \(S_3\)-sets or permutation modules.  There
is an abstract \(S_3\)-module isomorphism
\begin{equation}
 \psi:\Lambda_3\xrightarrow{\sim}\mathscr R,qquad
 \psi(e_i,0)=r_i^+,quad \psi(0,e_i)=r_i^-,             \label{eq:rank6-formal-ray-iso}
\end{equation}
but it is not induced by the fan: both geometric ray signs in
\eqref{eq:rank6-chamber} come from \(L_3\), while \(\bar L_3\) records deck
and filling data.

\subsection{Spectral-floor evaluation and the monodromy obstruction}

Let
\[
 \mathscr D_{\mathbb Z}=\bigoplus_{i=1}^3\mathbb ZE_{ii}
 \subset\operatorname{Herm}_3(\mathbb C),
 \qquad
 F(H)=H-\lambda_{\min}(H)I_3.
\]
The diagonal Weyl action is conjugation by permutation matrices.  Since
\[
 F(\nu_i)=E_{ii},\qquad F(-\nu_i)=I_3-E_{ii},
\]
the endpoint values define the unique \(S_3\)-equivariant homomorphism
\begin{equation}
 E:\mathscr R\longrightarrow\mathscr D_{\mathbb Z},qquad
 E(r_i^+)=E_{ii},\quad E(r_i^-)=I_3-E_{ii}.              \label{eq:rank6-ray-evaluation}
\end{equation}

\begin{proposition}[Exact diagonal Gram-boundary carrier]
\label{prop:rank6-gram-carrier}
The map \(E\) is surjective and
\begin{equation}
 \ker E=\mathbb Z\left\langle
 \begin{array}{l}
 r_1^--r_2^+-r_3^+,\\
 r_2^--r_1^+-r_3^+,\\
 r_3^--r_1^+-r_2^+
 \end{array}\right\rangle.                              \label{eq:rank6-e-kernel}
\end{equation}
For the trace pairing on diagonal matrices, the pullback matrix on
\(\mathscr R\) is
\begin{equation}
 E^{\mathsf T}E=
 \begin{pmatrix}
 I_3&J_3-I_3\\
 J_3-I_3&I_3+J_3
 \end{pmatrix},                                        \label{eq:rank6-e-pairing}
\end{equation}
where \(J_3\) is the all-ones matrix; its radical is \(\ker E\).
For every chamber,
\begin{equation}
 G_\pi:=(EI)(c_\pi)
 =2E_{\pi_1\pi_1}+E_{\pi_2\pi_2}
 =F(\nu_{\pi_1}-\nu_{\pi_3}).                         \label{eq:rank6-six-gram-points}
\end{equation}
The six \(G_\pi\) are distinct rank-two positive-semidefinite matrices and
form a free \(S_3\)-orbit on the diagonal part of the cubic Gram boundary.
Their integral span is
\begin{equation}
 \left\{\operatorname{diag}(m_1,m_2,m_3):
 m_1+m_2+m_3\equiv0\pmod3\right\},                     \label{eq:rank6-gram-span}
\end{equation}
an index-three sublattice of \(\mathscr D_{\mathbb Z}\).
\end{proposition}

\begin{proof}
The positive-ray images already form the standard basis, so \(E\) is
surjective.  In the displayed ray order its matrix is
\([I_3\mid J_3-I_3]\).  Solving its three coordinate equations gives
exactly the three independent generators in \eqref{eq:rank6-e-kernel};
multiplying the matrix by its transpose gives
\eqref{eq:rank6-e-pairing}.  On the chamber
\(s\nu_i+t(-\nu_j)\), where \(i\ne j\) and \(s,t\geq0\), the \(j\)-th
diagonal entry is least, and direct subtraction shows
\[
 F(s\nu_i-t\nu_j)=sE_{ii}+t(I_3-E_{jj}).
\]
Putting \(s=t=1\) proves \eqref{eq:rank6-six-gram-points}; the ordered
diagonal entries are a permutation of \((2,1,0)\), which proves rank, positivity,
distinctness, and freeness.  Every such vector has coordinate sum three.
Conversely, the three columns corresponding to
\(G_{123},G_{132},G_{213}\) have determinant \(-3\).  Their span therefore
has index three and lies in the index-three lattice on the right of
\eqref{eq:rank6-gram-span}; equality follows.
\end{proof}

The preceding proposition is a typed map from the chamber/ray module to a
rank-three diagonal lattice in the nine-real-dimensional Hermitian carrier.
It does not yet give a map from all of \(\Lambda_3\), and the missing
compatibility cannot be repaired by assigning the deck basis to the negative
rays.  Indeed, the homomorphism
\(g(e_i)=E_{ii}\) used below is the unique linear extension of the three
positive endpoint labels to \(L_3\); it is not a map on the quotient
\(A_F=L_3/\mathbb Zu\), since \(g(u)=I_3\ne0\).  The genuine spectral-floor
map is only piecewise linear on the chamber fan.

\begin{theorem}[No monodromy-compatible extension of the angular Gram rays]
\label{thm:rank6-monodromy-obstruction}
Let
\[
 g:L_3\xrightarrow{\sim}\mathscr D_{\mathbb Z},qquad
 g(a)=\operatorname{diag}(a_1,a_2,a_3).
\]
There do not exist homomorphisms
\(\eta:\Lambda_3\to\mathscr D_{\mathbb Z}\) and
\(T:\mathscr D_{\mathbb Z}\to\mathscr D_{\mathbb Z}\) such that
\begin{equation}
 \eta|_{L_3\oplus0}=g,qquad \eta M_3=T\eta.             \label{eq:rank6-impossible-descent}
\end{equation}
Consequently no torus homomorphism whose angular cocharacter map is the
Gram-ray marking can intertwine the displayed monodromy with an endomorphism
of this diagonal Gram carrier.
\end{theorem}

\begin{proof}
For every \(a\in L_3\), \(M_3(a,0)=(a,0)\).  Hence
\eqref{eq:rank6-impossible-descent} gives \(Tg(a)=g(a)\).  Since \(g\) is
surjective, \(T=I\).  For every \(b\in\bar L_3\), the same equation applied
to \((0,b)\) now gives
\[
 g(b)+\eta(0,b)=\eta(b,b)=\eta M_3(0,b)
 =T\eta(0,b)=\eta(0,b),
\]
so \(g(b)=0\) for all \(b\), contradicting that \(g\) is an isomorphism.
A torus homomorphism induces a homomorphism of cocharacter lattices, and
intertwining family monodromies induces the second equality in
\eqref{eq:rank6-impossible-descent}; therefore the lattice contradiction is
also the stated torus-level obstruction.
\end{proof}

For completeness, the tempting formal assignment is
\begin{equation}
 \eta_0=E\psi,qquad
 \eta_0(a,b)=\operatorname{diag}(a)
 +(u^{\mathsf T}b)I_3-\operatorname{diag}(b).            \label{eq:rank6-eta0}
\end{equation}
Its kernel has the basis
\begin{equation}
 k_i=(e_i-u,e_i),\qquad i=1,2,3.                         \label{eq:rank6-eta0-kernel}
\end{equation}
This kernel is a rank-three Lagrangian for \(\Omega_3\), because direct
substitution gives \(\Omega_3(k_i,k_j)=0\).  It is not monodromy invariant:
\begin{equation}
 \eta_0(k_i)=0,qquad \eta_0(M_3k_i)=E_{ii}.             \label{eq:rank6-eta0-failure}
\end{equation}
Thus \eqref{eq:rank6-eta0-failure} exhibits the obstruction of
\cref{thm:rank6-monodromy-obstruction} on three primitive generators.

\subsection{The branch map is angular: the corrected \texorpdfstring{\(\beta\)}{beta} square}
\label{subsec:rank6-beta-square}

Consider the monomial torus map
\[
 q:(\mathbb C^*)^4\longrightarrow(\mathbb C^*)^3,qquad
 q(z_0,z_1,z_2,z_3)=(z_0z_3,z_1,z_2).
\]
On cocharacter lattices it induces
\begin{equation}
 \beta:\mathbb Z^4\longrightarrow\mathbb Z^3,qquad
 \beta(n_0,n_1,n_2,n_3)=(n_0+n_3,n_1,n_2).             \label{eq:rank6-beta}
\end{equation}
If \(\Sigma_r\) denotes coordinate sum on \(\mathbb Z^r\), then
\begin{equation}
 \Sigma_3\beta=\Sigma_4.                                \label{eq:rank6-sum-square}
\end{equation}
Thus \(\beta\) restricts from the four-branch angular fibre lattice
\(\ker\Sigma_4\) to the three-branch angular fibre lattice
\(\ker\Sigma_3\).

Define the integral isomorphisms
\begin{align}
 \alpha_3:L_3&\xrightarrow{\sim}\ker\Sigma_4,&
 \alpha_3(x_1,x_2,x_3)
 &=(-x_1,x_1-x_2,x_2-x_3,x_3),                          \label{eq:rank6-alpha3}\\
 \alpha_2:A_F&\xrightarrow{\sim}\ker\Sigma_3,&
 \alpha_2(y_1,y_2)&=(-y_1-y_2,y_1,y_2).                 \label{eq:rank6-alpha2}
\end{align}

\begin{proposition}[Typed angular factorization]
\label{prop:rank6-beta-p}
The exact branch/fan square is
\begin{equation}
 \boxed{\ \beta\alpha_3=\alpha_2P\ },                 \label{eq:rank6-beta-p}
\end{equation}
or equivalently
\[
 \begin{CD}
 L_3 @>{P}>> A_F\\
 @V{\alpha_3}V{\cong}V @VV{\alpha_2}V\\
 \ker\Sigma_4 @>{\beta|_{\ker\Sigma_4}}>> \ker\Sigma_3.
 \end{CD}
\]
The deck map \(Q:\bar L_3\to B_F\) belongs instead to the separate filling
square \(c_F\Phi=Qc_3\).  Therefore the written matrix product \(Q\beta\),
although numerically formable after erasing all lattice labels, is not a
composition of the displayed geometric maps: the codomain of \(\beta\) is
the branch cocharacter lattice \(\mathbb Z^3\), whereas the domain of \(Q\)
is the independently marked deck lattice \(\bar L_3\).
\end{proposition}

\begin{proof}
The matrices of the two vertical maps in the displayed bases are
\[
 [\alpha_3]=
 \begin{pmatrix}-1&0&0\\1&-1&0\\0&1&-1\\0&0&1\end{pmatrix},
 \qquad
 [\alpha_2]=
 \begin{pmatrix}-1&-1\\1&0\\0&1\end{pmatrix}.
\]
They have images in the stated sum-zero lattices, and solving successively
for the partial sums gives integral inverses.  Multiplication with
\eqref{eq:rank6-pq} and \eqref{eq:rank6-beta} gives
\[
 [\beta][\alpha_3]
 =\begin{pmatrix}-1&0&1\\1&-1&0\\0&1&-1\end{pmatrix}
 =[\alpha_2]P,
\]
which proves the square.  The source and target labels in
\eqref{eq:rank6-phi} show that \(Q\) acts on \(\bar L_3\), not on the branch
coordinate lattice appearing in \eqref{eq:rank6-beta}; this proves the type
statement.  The valid deck/filling identity was already proved in
\cref{thm:rank6-exact-sequence}.
\end{proof}

\subsection{The distinct regular-locus defect lattice}
\label{subsec:distinct-defect-lattice}

For the candidate six-dimensional regular locus
\(U=X\setminus(W\sqcup S_1\sqcup S_2)\), the independent defect calculation
has the marked decomposition
\begin{equation}
 \mathcal L_{\rm def}:=H^4(U;\mathbb Z)
 \cong H_1(W;\mathbb Z)\oplus H_1(S_1;\mathbb Z)
               \oplus H_1(S_2;\mathbb Z)
 \cong(\mathbb Z^2)^3.                                  \label{eq:rank6-defect-lattice}
\end{equation}
The cusp, order-three, and order-four support components distinguish the
three summands; no geometric \(S_3\)-action permuting them is supplied.  In
the marked basis \((a_0,b_0,a_1,b_1,a_2,b_2)\), its primitive fibre-degree
map is
\begin{equation}
 \deg_F:\mathcal L_{\rm def}\longrightarrow\mathbb Z,
 \qquad [\deg_F]=(12,0,4,0,-3,0).                       \label{eq:rank6-defect-degree}
\end{equation}
Its reductions define finite character quotients such as
\(\mathcal L_{\rm def}/\ker(\rho_6\deg_F)\cong\mathbb Z/6\), but this is
secondary quotient data of \(\mathcal L_{\rm def}\), not its underlying
rank-six structure.

There is consequently no typed identification in the retained data between
\(\mathcal L_{\rm def}\) and \(\Lambda_3\).  Equality of ranks supplies no
map; their markings, functorial origins, distinguished maps, and available
pairings differ.  In particular \(\Omega_3\) is the explicitly chosen
nondegenerate alternating form in the marked decomposition
\(L_3\oplus\bar L_3\), not an intersection pairing canonically attached to
the six-torus.  Ordinary Poincar\'e duality for the noncompact
six-dimensional regular locus pairs \(H^4(U)\) with compactly supported
degree-two data and does not provide the displayed self-pairing on
\(\mathcal L_{\rm def}\).  A comparison would require an explicit geometric
correspondence
\[
 H_1(T^6;\mathbb Z)\longrightarrow H^4(U;\mathbb Z)
\]
together with proofs of its compatibility with monodromy, support markings,
and the relevant pairings.  No such correspondence occurs in the audited
construction.

\subsection{What the exact carrier does and does not supply}

The verified maps now form three distinct diagrams:
\[
 \Lambda_3\xrightarrow{P\oplus Q}A_F\oplus B_F,
 \qquad
 \mathscr C\xrightarrow{I}\mathscr R\xrightarrow{E}\mathscr D_{\mathbb Z},
 \qquad
 \ker\Sigma_4\xrightarrow{\beta}\ker\Sigma_3.
\]
The first intertwines monodromy and filling, the second records angular
Weyl-chamber incidence and its diagonal spectral-floor image, and the third
is the cocharacter map of the local monomial toric morphism.  The equality
of several ranks or the possibility of multiplying their unlabeled integer
matrices does not merge these diagrams.  In particular, the construction
provides six exact marked points in the rank-two stratum of the cubic Gram
boundary, but it does not provide a continuous, holomorphic, toric, or
stratified map from the compactified toroidal total space to the full
eight-dimensional Gram boundary.  Such a global map would still require
explicit overlap descent, a singular-fibre extension, and a boundary
comparison compatible with monodromy and filling; the lattice obstruction
of \cref{thm:rank6-monodromy-obstruction} rules out the most direct angular
Gram-ray extension.

\section*{Conclusion}

Two complex color triplets produce an exact four-real-dimensional Gram orbit
space, and the compact representation quotient is the four-ball with the
complete orbit-type boundary displayed above.  Independently, the toroidal
homology lattice \(\Lambda_3\) carries an exact diagonal Weyl action,
symplectic shear, and weight-lattice quotient.  Its six Kuhn chambers map by
an integral incidence complex to six diagonal rank-two points of the cubic
Gram boundary.  They do not identify with the six basis directions of
\(\Lambda_3\), and no lattice map extending the angular Gram-ray marking can
intertwine the toroidal monodromy.  The branch map factors through \(P\), not
through the deck map \(Q\), and the separately arising free defect lattice
\(H^4(U;\mathbb Z)\) remains nonidentified with \(\Lambda_3\).

Thus the surviving result is an exact collection of typed orbit, lattice,
fan, and boundary-carrier theorems together with a sharp monodromy
obstruction.  It does not supply the missing global toroidal-to-Gram map, a
nonzero \(S^4\) clutching class, or a nontrivial interacting Yang--Mills
sector.

\end{document}
